\documentclass[10pt,a4paper,ssfamily]{exam}
\usepackage[brazil]{babel}
\usepackage[numbers,sort&compress]{natbib}
\usepackage[utf8]{inputenc}
\renewcommand{\baselinestretch}{1.0}
\usepackage{epsfig}
\usepackage{mhchem}
\usepackage{graphicx}
\usepackage{makeidx}
\usepackage{multicol}
\usepackage{multirow}
\usepackage{amssymb}
\usepackage{amsmath}
\DeclareMathOperator{\sen}{sen}
\newcommand{\listfont}{\bf\fontencoding{OT1}\sffamily\large}
\newcommand{\titlesfont}{\fontencoding{OT1}\sffamily\large}
\renewcommand{\baselinestretch}{1.0}
\newcommand{\Mg}{Mg$^{2+}$}
\renewcommand{\t}{\textrm}
\newcommand{\cc}[1]{[\textrm{#1}]}
\usepackage{enumitem}
\newcommand{\bmat}{\left[\begin{array}{cc}}
\newcommand{\emat}{\end{array}\right]}
\newcommand{\nomera}{
  \noindent
  \begin{tabular}{|p{0.7\textwidth}|p{0.3\textwidth}|}
  Nome: & RA: \\
  \hline
  \end{tabular}
}

\newcommand{\qini}{\begin{enumerate}[resume]\item}
\newcommand{\qend}{\end{enumerate}}

\newcommand{\oC}{$^\circ$C}
\newcommand{\kcalmol}{~kcal~mol$^{-1}$}
\newcommand{\moll}{~mol~L$^{-1}$}
\newcommand{\molkg}{~mol~kg$^{-1}$}
\newcommand{\gl}{~g~L$^{-1}$}
\newcommand{\e}[1]{$\times 10^{#1}$}

\newcommand{\eps}{\varepsilon}

\newcommand{\Resposta}[1]{\vspace{0.5cm}\noindent{\bf Resposta:}\\ #1
\begin{center}\rule{\textwidth}{0.4pt}\end{center}}
\renewcommand{\Resposta}[1]{}
\newcommand{\resposta}[1]{\vspace{0.5cm}\noindent{\bf Resposta:}\\ #1
\begin{center}\rule{\textwidth}{0.4pt}\end{center}}


\begin{document}
\sffamily
\pagestyle{empty}
\nomera

\begin{center}

{\bf QG101 - Química Geral}

{\bf Aula 4 - Configuração Eletrônica e Tabela Periódica}\\
\underline{Justifique suas respostas.}
\end{center}

\begin{center}{\bf Questões}\end{center}

% ---------------------------------------------------------------
\qini
O \textbf{manganês} (\ce{Mn}, $Z = 25$) é amplamente utilizado em ligas
metálicas de aço e em pilhas alcalinas (\ce{MnO2}). Sua química variada
está diretamente ligada à sua configuração eletrônica.

\begin{enumerate}
  \item Escreva a configuração eletrônica completa do \ce{Mn} e sua
        forma abreviada usando o gás nobre anterior. Identifique a qual
        bloco, período e grupo da tabela periódica ele pertence.
  \item Represente o preenchimento dos orbitais 3d do \ce{Mn} usando
        diagramas de caixas ({\footnotesize $[\uparrow\downarrow]$,
        $[\uparrow\phantom{\downarrow}]$, $[\phantom{\uparrow}\phantom{\downarrow}]$}).
        Quantos elétrons desemparelhados possui? O \ce{Mn} é
        paramagnético ou diamagnético?
  \item O \ce{Mn} forma o íon \ce{Mn^{2+}} (perde 2 elétrons) e o íon
        \ce{Mn^{4+}} (perde 4 elétrons). Escreva a configuração
        eletrônica de cada íon e indique de qual(is) subnível(is) os
        elétrons foram removidos, justificando a ordem de remoção.
\end{enumerate}

\vspace{2cm}
\qend

% ---------------------------------------------------------------
\qini
O \textbf{silício} (\ce{Si}, $Z = 14$) e o \textbf{germânio}
(\ce{Ge}, $Z = 32$) são semicondutores fundamentais na indústria
eletrônica. Suas propriedades dependem diretamente da configuração
eletrônica de valência.

\begin{enumerate}
  \item Escreva a configuração eletrônica completa e abreviada do
        \ce{Si} e do \ce{Ge}. Quantos elétrons de valência cada um
        possui?
  \item Ambos pertencem ao mesmo grupo da tabela periódica. Qual grupo?
        Justifique com base na configuração de valência.
  \item Na dopagem tipo~n, átomos de fósforo (\ce{P}, $Z = 15$)
        substituem átomos de \ce{Si} na rede cristalina. Escreva a
        configuração de valência do \ce{P} e explique, em termos de
        elétrons de valência, por que essa substituição gera um
        elétron ``extra'' livre para condução.
\end{enumerate}

\vspace{2cm}
\qend

% ---------------------------------------------------------------
\qini
O \textbf{cromo} (\ce{Cr}, $Z = 24$) é utilizado em revestimentos
anticorrosivos e em ligas inoxidáveis. Sua configuração eletrônica
apresenta uma exceção em relação ao diagrama de Pauling.

\begin{enumerate}
  \item Qual seria a configuração eletrônica do \ce{Cr} prevista pelo
        diagrama de Pauling? E qual é a configuração real? Explique
        por que a configuração real é mais estável.
  \item O \ce{Cr} pode formar os íons \ce{Cr^{3+}} e \ce{Cr^{6+}}.
        Escreva a configuração eletrônica de cada um. Qual deles é
        diamagnético? Justifique.
  \item O cobre (\ce{Cu}, $Z = 29$) apresenta anomalia semelhante à
        do \ce{Cr}. Escreva a configuração real do \ce{Cu} e
        explique a estabilidade extra envolvida.
\end{enumerate}

\vspace{2cm}
\qend

\end{document}
